\subsection{Exact reconstruction in a specified harmonic neighbourhood}
\label{sec:harmonic-reconstruction}

We isolate the local theorem used to complete the construction.
Its low-order coefficients are explicit hypotheses, derived for our
special map in the preceding subsections.  The proof has two independent
parts: necessary equations have at most one solution, and a small exact
model supplies a solution.  This separation is what permits an
all-orders conclusion without extrapolating from finite precision.

Let $\mathcal O$ be the ring of integers of
$K_{0,23}=\QQ_{23}(\pi)$, where $\pi^2=-23$.
In canonical coordinates $[H:U:V:W]$ put
\[
 q_0=HW-U^2-UV-V^2,\qquad c_0=UW^2-V^3.
\]
The normalization of their common zero curve in characteristic $23$
has, on $W=1$, the parametrization
\[
 (H,U,V)=(z^6+z^4+z^2,z^3,z).
\]
The two marks are $z=1,-1$; the critical cluster is at $z=0$.
In these coordinates the normalized reduced map is
$(1-z^{23})/(1+z^{23})$.  Thus
$t_{\mathrm n}=(1-z)/(1+z)$ relates this description to
\cref{prop:harmonic-position}; the singular point $z=\infty$
has $t_{\mathrm n}=-1$.

Write
\[
 (m_0,\ldots,m_8)
 =(H^3,H^2U,HU^2,U^3,H^2V,HUV,U^2V,HV^2,UV^2)
\]
and set
\begin{align*}
 C_1&=H^3-5HU^2-7HUV+11HV^2,\\
 C_2&=6H^2U+5U^3+3H^2V+4U^2V+3UV^2,\\
 C_3&=-H^3-10HU^2+6HUV+2HV^2.
\end{align*}
The neighbourhood considered below consists of equations and marks
of the form
\begin{equation}\label{eq:harmonic-neighbourhood}
\begin{aligned}
 q&=q_0,\\
 c&=c_0+\pi C_1+\pi^2C_2+\pi^3C_3
       +\pi^3 u_0UV^2+\pi^4\sum_{i=0}^{7}u_{i+1}m_i,\\
 z_A&=1+\pi+9\pi^2+\pi^3u_9,\\
 z_B&=-1+\pi-9\pi^2+\pi^3u_{10},
 \qquad u_i\in\mathcal O.
\end{aligned}
\end{equation}
The remaining coordinates of the marks are determined on the smooth
central chart by $q=c=0$.

\begin{theorem}[Incidence in the specified local neighbourhood]
\label{thm:harmonic-local-incidence}
Suppose a smooth genus-$4$ curve and a degree-$23$ three-point map
admit the canonical equations and ordered totally ramified marks in
\cref{eq:harmonic-neighbourhood}.  Suppose its remaining ramification
consists of eight simple points in the formal neighbourhood of $z=0$,
and that the integral differential
\[
 \eta=\frac{1}{23}\frac{d\beta}{\beta}
 \quad\hbox{has reduction}\quad
 \overline\eta=\frac{2z^8}{z^2-1}\,dz,
 \qquad \operatorname{div}(\beta)=23A-23B.
\]
Then the marked canonical curve is unique up to the coordinate
normalizations just specified.  It is isomorphic over $K_{0,23}$
to the exact model below.  Its two osculating planes meet the curve
in two distinct geometric points.

The same conclusion holds over an unramified complete extension of
$\mathcal O$ with the same congruences.  The exact model has a
degree-$23$ map with geometric monodromy $M_{23}$ and represents the
distinguished member of the seven-cover family.
\end{theorem}

\paragraph{An intrinsic choice of coordinates.}
For $i=0,1,2,3$, the spaces
\[
 H^0(X,\omega_X(-iA-(3-i)B))
\]
are one-dimensional in this neighbourhood, and their generators form
a basis.  This follows from rank-three jet minors and a nonzero
basis determinant on the closed fibre.  Use these generators as
coordinates $x_0,x_1,x_2,x_3$.  The marks become
$[1:0:0:0]$ and $[0:0:0:1]$, with vanishing sequences
$(0,1,2,3)$ and $(3,2,1,0)$.
The quadric consequently has the form
\[
 ax_0x_3+bx_0x_2+dx_1x_2+ex_1^2+fx_1x_3+gx_2^2.
\]
Here $a,b,d,e$ are units.  Substituting $x_i=d_i y_i$ and multiplying
the equation by $\kappa$, with
\[
 d_0=1,\qquad d_1=b/d,\qquad
 \kappa=(ed_1^2)^{-1},\qquad
 d_2=(\kappa b)^{-1},\qquad d_3=(\kappa a)^{-1},
\]
sets the first four displayed coefficients to one.
This normalization is unique and requires no extension of the field.
Reduce the cubic modulo the monic quadric in lexicographic order
$x_0>x_1>x_2>x_3$, and set its $x_0^2x_3$ coefficient to one.
These operations remove the scalar choices in the four generators.

\paragraph{The exact model.}
In these normalized coordinates the quadric is
\begin{equation}\label{eq:small-harmonic-quadric}
 q_*=x_0x_2+x_0x_3+x_1^2+x_1x_2
             -\frac{127}{229}x_1x_3+\frac{12192}{52441}x_2^2.
\end{equation}
The cubic $c_*$ has the following nonzero coefficients; all others
vanish.
\begin{center}
\renewcommand{\arraystretch}{1.18}
\begin{tabular}{@{}ll@{}}
\toprule
Monomial & Coefficient in $c_*$\\
\midrule
$x_0^2x_3$ & $1$\\
$x_0x_1x_3$ & $(69+41\pi)/916$\\
$x_0x_3^2$ & $-(288+3552\pi)/52441$\\
$x_1^3$ & $-(5+\pi)/458$\\
$x_1^2x_2$ & $-(3061+569\pi)/209764$\\
$x_1^2x_3$ & $-(8649+1515\pi)/104882$\\
$x_1x_2^2$ & $-(298479+49803\pi)/48035956$\\
$x_1x_2x_3$ & $-(1041970+181274\pi)/12008989$\\
$x_1x_3^2$ & $(2318419+389759\pi)/48035956$\\
$x_2^3$ & $-(2350008+338328\pi)/2750058481$\\
$x_2^2x_3$ & $-(55642056+9354216\pi)/2750058481$\\
\bottomrule
\end{tabular}
\end{center}

\begin{proof}[Proof of \cref{thm:harmonic-local-incidence}]
We separate local uniqueness, exact existence, and identification of
the cover.  In particular, recognizing coefficients from a local
approximation is not used as a proof of their exact values.

\emph{Local uniqueness.}
The nine cubic directions $m_i$ give a formal slice for canonical
curves with quadric $q_0$: the six infinitesimal orthogonal-coordinate
directions, cubic scaling, the four multiples of $q_0$, and these
nine directions span the twenty-dimensional cubic space.
The determinant is $10$ in $\FF_{23}$ for the certificate's basis.
The formal inverse-function theorem therefore justifies this slice.

Choose bases $S_A,S_B$ of the four-dimensional kernels of the
order-$23$ quintic jet maps.  Fixed unit minors make their
coefficients analytic functions of the eleven parameters $u_i$.
An invertible minor of fifteen cross-product equations determines
a multiplier matrix $M$, with one entry fixed to remove its scalar.
It defines a local ratio $N/D$, normalized to value one at $z=0$,
even off the locus where all global cross-products vanish.
On a genuine cover this is its function $\beta$, up to scalar.
No global map is inferred merely from solving those fifteen rows.

Let $R(z)$ be the degree-eight Weierstrass factor of $\eta$, with
$R\equiv z^8\pmod\pi$.  Reduce $N-D$ and $D$ modulo $R^2$, obtaining
$P,Q$, and put
\[
 E=P-\frac{P(0)}{Q(0)}Q,\qquad E_j=[z^j]E.
\]
Here $Q(0)$ is a unit.  Retaining $R^2$ retains the effective
ramification multiplicities.  For a three-point map $E=0$,
because its eight simple critical points have one common value.
Let $C_j$ denote the central-series coefficient of $z^j$ in the
negative of the first remaining cross-product
$S_{A,0}(MS_B)_1-S_{A,1}(MS_B)_0$, using the fixed bases in the
certificate.  Eleven necessary equations are
\begin{equation}\label{eq:harmonic-hensel-equations}
 F=\left(\frac{E_{15}}{\pi^3},
       \frac{E_9}{\pi^4},\ldots,\frac{E_{14}}{\pi^4},
       \frac{C_7}{\pi^5},\ldots,\frac{C_{10}}{\pi^5}\right)=0.
\end{equation}
These are integral restricted power series on the stated neighbourhood.
The exact residual calculation gives
\[
 \overline F=J\overline u+b,\qquad
 \det J=16\in\FF_{23}^{\times},\qquad
 -J^{-1}b=(0,0,-7,0,-1,-2,0,10,0,-7,-7).
\]
The full matrix, fixed jet bases, and normalizations are retained in
the local certificate and its accompanying derivation.

This calculation is an identity test, not a sample of possible covers.
The parameters first enter at orders $\pi^3,\pi^4$.
The first variation of the normalized ratio vanishes modulo $\pi$;
its variable contribution begins at order $\pi^4$.
After division by $23=-\pi^2$, the coefficients of $R$ first vary
at order $\pi^2$.  Modulo the precisions needed in
\cref{eq:harmonic-hensel-equations}, all expressions have parameter
degree at most two.  Values at zero, the positive and negative
coordinate vectors, and the sums of two distinct coordinate vectors
therefore determine them.  The $78$ evaluations verify integrality
of the divisions and the displayed affine reduction.
The computation also verifies that the two unavailable final digits
of the Weierstrass factor do not affect these divided reductions.

Consequently $F(u)=\widetilde J u+\widetilde b+\pi G(u)$, with
$G$ integral and restricted and $\widetilde J$ invertible over
$\mathcal O$.  The map
$u\mapsto-\widetilde J^{-1}\widetilde b
          -\pi\widetilde J^{-1}G(u)$
is a strict contraction on $\mathcal O^{11}$.
It has a unique fixed point, also after the stated base extensions.
Every cover in the theorem satisfies these necessary equations.
We do not need to assert that the selected equations alone imply
all the other map equations.

\emph{Exact existence.}
The pair $(q_*,c_*)$ was recovered from the local lift in the
intrinsic coordinates above, without using the previously stored
generic models.  The following independent tests are over $K_0$.
The Jacobian-minor ideal is the unit ideal on all four affine
projective charts, so the curve is a smooth genus-$4$ complete
intersection.  Its marked vanishing sequences are as stated.
The order-$23$ quintic jet kernels each have dimension four.
Their common base divisors are exactly $23A$ and $23B$: saturation
excludes other base points, and a section has order exactly $23$
at each mark.

In the degree-ten canonical ring, of dimension $57$, an exact
rank-$15$ multiplier calculation and all six cross-product
identities give a function with divisor $23A-23B$.
Away from its zero and pole, a critical ideal in the chart $x_1=1$
has an eight-dimensional reduced quotient algebra, and $N=uD$
in that algebra for one nonzero scalar $u$.
Reducedness is certified by a squarefree degree-eight characteristic
polynomial of a multiplication operator.
Riemann--Hurwitz gives total ramification $52$; the two marked
points contribute $44$, and the eight distinct critical points
exhaust the remainder.  Thus each is simple, there are no further
critical points in another chart, and the map has the required
three-point passport.

An integral projective change with unit determinant identifies
this exact marked model with the stated local slice modulo
$\pi^{33}$.  Correcting the quadric and cubic by the formal slice
construction preserves the lower-order congruences in
\cref{eq:harmonic-neighbourhood}.  The normalized map is determined
by its divisor; its divided differential and critical cluster are
therefore those used to form the necessary equations above.
The exact map belongs to the required neighbourhood and hence,
by uniqueness, is the local lift to all orders.

\emph{Incidence and identification.}
The two osculating planes on this exact model are $x_3=0$ and
$x_0=0$.  On their common line the identity
\begin{align*}
 c_*(0,x_1,x_2,0)
 &=\left(x_1^2+x_1x_2+\frac{12192}{52441}x_2^2\right)\\
 &\hspace{5mm}\cdot
 \left(-\frac{5+\pi}{458}x_1
             -\frac{771+111\pi}{209764}x_2\right)
\end{align*}
has first factor equal to $q_*|_{x_0=x_3=0}$.
Its discriminant is $3673/52441\ne0$.  This proves the claimed
two reduced points, exactly rather than to finite precision.

Finally, a separate comparison certificate constructs an invertible
projective matrix between this model and the degree-one model of
\cref{sec:hurwitz-scheme}.  It sends the ordered marks to the
ordered marks and verifies equality of the two rational functions
up to a nonzero target scalar by exact reduction in the canonical
ring.  The existing certified branch cycles consequently give
geometric monodromy $M_{23}$ for this map.  This is an exact
isomorphism argument, not an inference from the passport and not
a new numerical continuation calculation.
\end{proof}

\Needspace{8\baselineskip}
\begin{remark}[An explicit descent of the reconstructed map]
The exact model also admits a direct descent, independently of
the exclusion of incidence on the other six covers.
Put $a=-127/229$, $b=12192/52441$, and define
\[
 T(x_0,x_1,x_2,x_3)=(abx_3,bx_2,x_1,x_0/a).
\]
If $\sigma(\pi)=-\pi$, exact substitution gives
\[
 T^2=bI,\qquad q_*(Tx)=bq_*(x),\qquad
 \sigma(c_*)(Tx)\equiv\mu c_*(x)\pmod{q_*},
\]
where $\mu=-(52130760+33951768\pi)/1525141603$ and
$\mu\sigma(\mu)=b^3$.
Thus $R=(b/\mu)T$ satisfies $R\sigma(R)=I$.
Solving $v=R\sigma(v)$ gives a four-dimensional $\QQ$-space
and an invertible change of coordinates in which the curve
equations are rational; the certificate checks these equations
directly.

For the map $f$ normalized to have third branch value $1$, one also
has $f(x)\sigma(f)(Tx)=1$ in the function field.
Consequently $t=(f-1)/(\pi(f+1))$ is fixed by the descent and
defines a map over $\QQ$.  Its branch locus is $t=0$ and
$23t^2+1=0$.  Rescaling by $T_{\mathrm{base}}=23t$ gives the
branch locus $0,\pm\sqrt{-23}$ used in the paper.
This recovers the known rational cover in different equations.
\end{remark}

\begin{remark}[Why uniqueness is enough]
The eleven equations are necessary for a cover in this neighbourhood.
We do not assert that their ideal is the entire map ideal.  Nor do
the displayed low-order incidence cancellations prove an exact
incidence relation on their own: an off-locus audit detects a higher
defect in the constant residual multiplier.
What closes the proof is that the constructed lift and the
independently verified exact model both satisfy the necessary system,
which has at most one solution.  Their identification transfers the
exact incidence calculation to the lift.  The earlier coefficient
checks place the model in the neighbourhood; uniqueness, not the
number of matching digits, supplies equality to all orders.
\end{remark}
